% GENERATED FILE. Edit canonical lessons, then run npm run generate:books.
% canonical-source-hash: fnv1a64:7d6f453ee90fc200
% canonical-lessons: ES-C730-encantador, ES-C730-encanto, ES-C730-atractivo, ES-C730-belleza, ES-C730-elegancia

\chapter{Encanto y belleza}
\label{ch:es-encanto-y-belleza}
\hlchaptermodality{730}

\begin{chapteropening}
\textbf{By the end of this chapter:} \emph{I can describe someone charming, attractive or elegant.}

Complete the last lesson of chapter 730: I can describe someone charming, attractive or elegant.
\end{chapteropening}

\section[encantador]{encantador — charming, delightful}

\label{lesson:ES-C730-encantador}



\begin{quote}
Before the new one: say the Spanish for nosy, then the Spanish for curiosity.
\end{quote}

\subsection*{You'll want to know: encantador}
\textbf{encantador} — \textquotedblleft{}charming, delightful\textquotedblright{}.

\begin{quote}
Tu madre es encantadora.
\end{quote}

\begin{quote}
\emph{Your mother is charming.}
\end{quote}

Feminine \emph{encantadora}. \emph{Un encantador de serpientes} is a snake charmer.

\begin{cousinweb}
\emph{encantador} is built on \emph{encantar}, which comes from Latin \emph{incantāre} \textquotedblleft{}to chant a spell over\textquotedblright{} (EtymDB), from \emph{cantāre} \textquotedblleft{}to sing\textquotedblright{}.
\end{cousinweb}

\begin{grammarlens}[title={What you've built}]
One more word for telling, explaining and describing at B1.
\end{grammarlens}

\subsection*{Your turn}
\begin{itemize}
\raggedright
  \item \emph{Say it:} \emph{encantador}
  \item \emph{Say it:} \emph{encantador}, once more
  \item \emph{Say it:} a sentence of your own with \emph{encantador}
  \item \emph{Recall:} say the Spanish for nosy, then the Spanish for curiosity, then say \emph{encantador} again
\end{itemize}

\begin{tcolorbox}[breakable,colback=teal!4,colframe=teal!35!black,title={Before you move on}]
What does encantador mean? (\textquotedblleft{}Charming, delightful\textquotedblright{}.) Say it once more.
\end{tcolorbox}
\section[el encanto]{el encanto — charm}

\label{lesson:ES-C730-encanto}



\begin{quote}
Before the new one: say the Spanish for curiosity, then the Spanish for charming.
\end{quote}

\subsection*{You'll want to know: el encanto}
\textbf{el encanto} — \textquotedblleft{}charm\textquotedblright{}.

\begin{quote}
El pueblo tiene mucho encanto.
\end{quote}

\begin{quote}
\emph{The village has a lot of charm.}
\end{quote}

\emph{Es un encanto} — he or she is a sweetheart.

\begin{cousinweb}
\emph{encanto} is built on \emph{encantar} \textquotedblleft{}to enchant\textquotedblright{}, from Latin \emph{incantāre} (EtymDB); charm was first a spell.
\end{cousinweb}

\begin{grammarlens}[title={What you've built}]
One more word for telling, explaining and describing at B1.
\end{grammarlens}

\subsection*{Your turn}
\begin{itemize}
\raggedright
  \item \emph{Say it:} \emph{el encanto}
  \item \emph{Say it:} \emph{el encanto}, once more
  \item \emph{Say it:} a sentence of your own with \emph{el encanto}
  \item \emph{Recall:} say the Spanish for curiosity, then the Spanish for charming, then say \emph{el encanto} again
\end{itemize}

\begin{tcolorbox}[breakable,colback=teal!4,colframe=teal!35!black,title={Before you move on}]
What does el encanto mean? (\textquotedblleft{}Charm\textquotedblright{}.) Say it once more.
\end{tcolorbox}
\section[atractivo]{atractivo — attractive}

\label{lesson:ES-C730-atractivo}



\begin{quote}
Before the new one: say the Spanish for charming, then the Spanish for charm.
\end{quote}

\subsection*{You'll want to know: atractivo}
\textbf{atractivo} — \textquotedblleft{}attractive\textquotedblright{}.

\begin{quote}
Es una oferta muy atractiva.
\end{quote}

\begin{quote}
\emph{It is a very attractive offer.}
\end{quote}

Also a noun: \emph{el atractivo de la ciudad} — the city's appeal.

\begin{cousinweb}
\emph{atractivo} comes from Latin \emph{attractīvus} (EtymDB), built on \emph{attrahere} \textquotedblleft{}to draw towards\textquotedblright{}, from \emph{trahere} \textquotedblleft{}to pull\textquotedblright{}.
\end{cousinweb}

\begin{grammarlens}[title={What you've built}]
One more word for telling, explaining and describing at B1.
\end{grammarlens}

\subsection*{Your turn}
\begin{itemize}
\raggedright
  \item \emph{Say it:} \emph{atractivo}
  \item \emph{Say it:} \emph{atractivo}, once more
  \item \emph{Say it:} a sentence of your own with \emph{atractivo}
  \item \emph{Recall:} say the Spanish for charming, then the Spanish for charm, then say \emph{atractivo} again
\end{itemize}

\begin{tcolorbox}[breakable,colback=teal!4,colframe=teal!35!black,title={Before you move on}]
What does atractivo mean? (\textquotedblleft{}Attractive\textquotedblright{}.) Say it once more.
\end{tcolorbox}
\section[la belleza]{la belleza — beauty}

\label{lesson:ES-C730-belleza}



\begin{quote}
Before the new one: say the Spanish for charm, then the Spanish for attractive.
\end{quote}

\subsection*{You'll want to know: la belleza}
\textbf{la belleza} — \textquotedblleft{}beauty\textquotedblright{}.

\begin{quote}
Me impresionó la belleza del paisaje.
\end{quote}

\begin{quote}
\emph{The beauty of the landscape struck me.}
\end{quote}

\emph{Un salón de belleza} — a beauty salon.

\begin{cousinweb}
\emph{belleza} is built on \emph{bello} \textquotedblleft{}beautiful\textquotedblright{} with the ending \emph{-eza} (EtymDB); \emph{bello} comes from Latin \emph{bellus} \textquotedblleft{}pretty, handsome\textquotedblright{}. Some dictionaries think the form came in through Catalan \emph{bellesa} or Occitan.
\end{cousinweb}

\begin{grammarlens}[title={What you've built}]
One more word for telling, explaining and describing at B1.
\end{grammarlens}

\subsection*{Your turn}
\begin{itemize}
\raggedright
  \item \emph{Say it:} \emph{la belleza}
  \item \emph{Say it:} \emph{la belleza}, once more
  \item \emph{Say it:} a sentence of your own with \emph{la belleza}
  \item \emph{Recall:} say the Spanish for charm, then the Spanish for attractive, then say \emph{la belleza} again
\end{itemize}

\begin{tcolorbox}[breakable,colback=teal!4,colframe=teal!35!black,title={Before you move on}]
What does la belleza mean? (\textquotedblleft{}Beauty\textquotedblright{}.) Say it once more.
\end{tcolorbox}
\section[la elegancia]{la elegancia — elegance}

\label{lesson:ES-C730-elegancia}



\begin{quote}
Before the new one: say the Spanish for attractive, then the Spanish for beauty.
\end{quote}

\subsection*{You'll want to know: la elegancia}
\textbf{la elegancia} — \textquotedblleft{}elegance\textquotedblright{}.

\begin{quote}
Viste con mucha elegancia.
\end{quote}

\begin{quote}
\emph{She dresses with great elegance.}
\end{quote}

The adjective is \emph{elegante}.

\begin{cousinweb}
\emph{elegancia} comes from Latin \emph{ēlegantia} \textquotedblleft{}refinement, elegance\textquotedblright{} (EtymDB), linked with \emph{ēligere} \textquotedblleft{}to pick out\textquotedblright{}: elegance is good choosing.
\end{cousinweb}

\begin{grammarlens}[title={What you've built}]
One more word for telling, explaining and describing at B1.
\end{grammarlens}

\subsection*{Your turn}
\begin{itemize}
\raggedright
  \item \emph{Say it:} \emph{la elegancia}
  \item \emph{Say it:} \emph{la elegancia}, once more
  \item \emph{Say it:} a sentence of your own with \emph{la elegancia}
  \item \emph{Recall:} say the Spanish for attractive, then the Spanish for beauty, then say \emph{la elegancia} again
\end{itemize}

\begin{tcolorbox}[breakable,colback=teal!4,colframe=teal!35!black,title={Before you move on}]
What does la elegancia mean? (\textquotedblleft{}Elegance\textquotedblright{}.) Say it once more.
\end{tcolorbox}
